From the figures of the Moon's phase depicted above, we can conclude that the
full moon occurs on June 11th ($\pm 0.5$ day), 2006. We are going to use this
full moon as the reference for calculating the Moon's age on Albert Einstein's
birthday (March 14, 1879).

The student must find out that the approximate time span between June 11th,
2006 and March 14, 1879 is 127 years.

There are 31 leap years between the two years (i.e.\ 1880, 1884, 1888, 1892,
1896, 1904, 1908, 1912, 1916, 1920, 1924, 1928, 1932, 1936, 1940, 1944, 1948,
1952, 1956, 1960, 1964, 1968, 1972, 1976, 1980, 1984, 1988, 1992, 1996, 2000,
and 2004).

So, there are $365 \times 127 + 31 = 46{,}386$ days starting from January 1,
2006 to January 1, 1879. As a consequence, there are 46,475 days starting from
June 11, 2006 to March 14, 1879.

Number of days / synodic month $= 46{,}475/29.5306 = 1573.79$
% note: the source prints "1573.791254" with the trailing digits "1254"
% struck through in red.
or 1573 synodic months and $0.791254 \times 29.5306 = 23.3662$ days.

The remainder of days tells us the number of days into the synodic cycle we've
gone. Since we are going backwards in time in these calculations (from 2006 to
1879), this is the number of days from the top of the cycle, rather than from
the bottom. To convert to the number of days from the bottom of the cycle
(full moon), we simply subtract this value from a full cycle (29.5306 days).
This is $29.5306 - 23.3662 = 6.1644$ days.

6.1644 days from the bottom of the cycle (full moon $= 14.7653$ days) or
$14.7653 + 6.1644 = 20.9297$ days into the synodic month. And that's our
answer.

Finally, the Moon's age on Albert Einstein's birthday (March 14, 1879) is
$\mathbf{20.9 \pm 0.5}$ \textbf{days}.

% FIGURE NEEDED: photograph of the waning gibbous Moon illustrating the
% appearance at age ~20.9 days; 2007 data-analysis solutions, page 3
% (question 2).
